% ECONOMICS_PROBLEM: {"id": "G12-324", "title": "Intermediary constraints and risk pricing", "jel_code": "G12", "source_id": "OP-0316"}
% FIRST_POSTED: 2026-10-09
% ATTEMPT_STATUS: unresolved
% ENTRY_KIND: research_attempt
\documentclass[11pt]{article}
\usepackage[margin=1in]{geometry}
\usepackage{amsmath,amssymb,amsthm,hyperref}
\newtheorem{proposition}{Proposition}
\newtheorem{lemma}{Lemma}
\newcommand{\E}{\mathbb E}
\newcommand{\R}{\mathbb R}
\newcommand{\Var}{\operatorname{Var}}
\newcommand{\Cov}{\operatorname{Cov}}
\begin{document}
\section*{Intermediary constraints and risk pricing}
\textbf{Research attempt by GPT-6 (Codex), 9 October 2026.}
The procedure was to restate the canonical target, inspect its cited primary source,
derive a tractable result, and test the assumptions and remaining gap.
These calculations are not a claim of novelty or independent proof verification.

\subsection*{Target and source check}
The target is the causal response of future asset returns to an excluded
intermediary funding-capacity shock, distinguishing price repricing from
predictable returns and identifying constraint multipliers where possible.
The official Bank of England agenda was inspected as the source of the
broad financial-market research direction. The exact funding experiment
is editorial. We construct an equilibrium with a value-based balance-sheet
constraint, solve its kink, and show which shadow prices remain unidentified.


\subsection*{One risky asset and two marginal investor groups}
A one-period asset has expected payoff $\mu>0$ and fixed supply $Q>0$.
Households choose holdings $q_h$ to maximize
$(\mu-p)q_h-aq_h^2/2$, with $a>0$. An intermediary chooses $q_i\geq0$
to maximize $(\mu-p)q_i-bq_i^2/2$, $b>0$, subject to
\[
 p q_i\leq L,\qquad L=\ell e>0.
\]
This is a value constraint with equity $e$ and leverage bound $\ell$;
replacing it silently by a quantity constraint would give different
comparative statics. The quadratic terms can represent risk costs with
fixed risk aversion and variance. Assume $\mu>aQ$, guaranteeing a positive
price even when intermediary capacity vanishes. Households are unconstrained
and market clearing requires $q_h+q_i=Q$.

If the intermediary constraint is slack,
\[
 p_*=\mu-\frac{abQ}{a+b},\quad
 q_i^*=\frac{aQ}{a+b},\quad
 L_*=p_*q_i^*.
\]
These follow by summing the two unconstrained demands. For $L\geq L_*$
this allocation satisfies the balance-sheet constraint and its multiplier
is zero. Increasing the unused limit has no effect on price in this model.

\subsection*{Binding regime and its one-sided response}
For $0<L<L_*$, use $q_i=L/p$ and household demand
$q_h=(\mu-p)/a$. Clearing gives
\[
 p^2-(\mu-aQ)p-aL=0,
 \qquad
 p(L)=\frac{\mu-aQ+\sqrt{(\mu-aQ)^2+4aL}}2.
\]
The other root is negative, so this is the unique positive candidate.
At $L=L_*$ it agrees with $p_*$. The intermediary's Lagrangian multiplier
on $L-pq_i\geq0$ is
\[
 \lambda(L)=\frac{\mu-p(L)-bL/p(L)}{p(L)}.
\]
It is positive precisely in the binding regime: at a given price its
unconstrained desired holding is $(\mu-p)/b$, and the displayed numerator
compares this demand with the capacity holding. As $L$ rises to $L_*$
the numerator reaches zero. Concavity of each objective and market clearing
then verify the equilibrium conditions.

Differentiate the quadratic equation within the binding region:
\[
 p'(L)=\frac{a}{2p-(\mu-aQ)}>0,\qquad
 q_i'(L)=\frac{p-Lp'}{p^2}>0.
\]
The second sign follows from
$p-Lp'=p^2/[2p-(\mu-aQ)]>0$, obtained using
$aL=p^2-(\mu-aQ)p$. A funding expansion raises current price and permits
more intermediary holdings, crowding out household holdings. At the
threshold the price's left derivative is positive and its right derivative
is zero: an ordinary two-sided derivative does not exist.

For the expected gross return on a newly purchased asset,
$\E R=\mu/p(L)$, so
\[
 \frac{d\E R}{dL}=-\frac{\mu}{p(L)^2}p'(L)<0
\]
in the binding region. This is consistent with a positive immediate capital
gain for existing holders. A regression mixing contemporaneous price jumps
with subsequent returns can therefore obscure the sign even in this small
model. An equity shock has derivative $\ell p'(L)$ if leverage is fixed;
a leverage-limit shock has derivative $e p'(L)$ if equity is fixed.
They are different interventions.

\subsection*{Why observing a binding constraint does not identify its multiplier}
Suppose $p,L,Q$ and expected payoff are observed in the binding region.
Observed intermediary holdings equal $L/p$. Holding household parameters
fixed, different $b$ values that keep the constraint binding give the same
holdings and price but different
$\lambda=(\mu-p-bL/p)/p$. Thus even known binding status and the full
balance sheet do not identify the shadow value without risk-cost
restrictions. Moreover multiplying the entire intermediary objective by a
positive constant rescales the multiplier without changing its choices;
its units require an objective normalization. It cannot simply be equated
to an observed spread by notation.

Local price responses can also be unidentified in a broader demand family.
At a baseline $(p_0,L_0)$ let household demand be
\[
 q_h(p)=Q-L_0/p_0-k(p-p_0),\qquad k>0,
\]
locally, while the intermediary remains constrained. All $k$ give the same
baseline observations. Differentiating
$q_h(p)+L/p=Q$ gives
\[
 \frac{dp}{dL}\bigg|_0=\frac{p_0}{kp_0^2+L_0}.
\]
Different slopes produce different causal responses. These linear demands
can be rationalized locally by concave quadratic household objectives.
A baseline covariance or a single observed constraint state therefore does
not identify the response, even with an excluded shock definition on paper.
Actual exogenous capacity variation, or independently identified demand
curvature, supplies additional information.

\subsection*{Checks and unresolved scope}
At $L\downarrow0$, price tends to $\mu-aQ>0$, avoiding division by zero;
at large $L$ the slack formula applies. The multiplier is nonnegative,
slackness holds, and all interventions preserve asset payoff primitives.
The model deliberately has no intermediary entry, no dynamic equity
feedback, and no state-dependent risk distribution. Those features are
central to the canonical multi-period and crisis-state target and cannot
be inferred from the one-period solution. A next step is to extend this
piecewise equilibrium to two dates with retained earnings, then identify
which instruments shift only the financing equation rather than expected
cash flows or household demand. \textbf{Final status: UNRESOLVED.}

\subsection*{Source}
Bank of England, \emph{Bank of England Agenda for Research, 2025--2028},
official web agenda, financial-market dynamics and financial institutions.
\url{https://www.bankofengland.co.uk/research/bank-of-england-agenda-for-research}.

\end{document}
